\section{Online normalization and minimal delay}\label{app:online-normalization}

This appendix records an explicit \emph{online} realization of Zeckendorf normalization and proves that a delay of $3$ is intrinsic.
The material is included as an auditable interface: it provides a constant-memory streaming implementation of the normal form that is compatible with the value-preserving rewrite system (Appendix~\ref{app:rewrite-system}).

\subsection{Extended normalization on digit streams}\label{app:online-norm-setup}
Recall the Fibonacci weights $(F_k)_{k\ge 1}$ with $F_1=F_2=1$.
For a finite-support digit configuration $a=(a_1,a_2,\dots)\in\mathcal{D}$ (Section~\ref{subsec:fold-rewrite}) define
\[
\mathrm{Val}(a):=\sum_{k\ge 1} a_k F_{k+1}.
\]
The rewrite system from Definition~\ref{def:fold-rewrite-rules} is strongly terminating and confluent on $\mathcal{D}$, hence every $a\in\mathcal{D}$ has a unique normal form $\mathrm{NF}(a)\in X_\infty$ with the same value (Theorem~\ref{thm:rewrite-unique-normal-form}).

\begin{definition}[One-sided normalization map]\label{def:fold-infty}
Define the one-sided normalization
\[
\Fold_\infty:\mathcal{D}\to X_\infty,\qquad \Fold_\infty(a):=\mathrm{NF}(a).
\]
For a finite word $s=s_1\cdots s_m\in\{0,1,2\}^m$, write $\iota_m(s)=(s_1,\dots,s_m,0,0,\dots)\in\mathcal{D}$ and set $\Fold_\infty(s):=\Fold_\infty(\iota_m(s))$.
\end{definition}

\subsection{A deterministic delay-$3$ online transducer}\label{app:online-delay-3-transducer}
Online normalization in Fibonacci numeration is classical in automata/numeration theory; see, e.g., \cite{Frougny1992NumbersAutomata,Frougny1999OnlineAddition}.
We record an explicit deterministic \emph{subsequential} transducer that realizes $\Fold_\infty$ on $\{0,1,2\}^\NN$ with delay $3$.

\begin{definition}[A $10$-state subsequential transducer]\label{def:delay3-transducer}
Let $\overline{1}:=-1$.
Consider the finite state set
\[
Q=\{000,001,002,010,100,101,0\overline{1}2,1\overline{1}2,01\overline{1},11\overline{1}\},
\]
with initial state $000$.
For each state $q\in Q$ and input digit $d\in\{0,1,2\}$, define a transition $q\xrightarrow{d/e}q'$ with output $e\in\{0,1\}$ as follows:
\[
\begin{aligned}
000&\xrightarrow{d/0}00d,\qquad 100\xrightarrow{d/1}00d\quad(d=0,1,2),\\
001&\xrightarrow{0/0}010,\ \xrightarrow{1/0}100,\ \xrightarrow{2/0}101,\\
002&\xrightarrow{0/0}11\overline{1},\ \xrightarrow{1/1}000,\ \xrightarrow{2/1}001,\\
010&\xrightarrow{0/0}100,\ \xrightarrow{1/0}101,\ \xrightarrow{2/1}0\overline{1}2,\\
101&\xrightarrow{0/1}010,\ \xrightarrow{1/1}100,\ \xrightarrow{2/1}101,\\
0\overline{1}2&\xrightarrow{0/0}01\overline{1},\ \xrightarrow{1/0}010,\ \xrightarrow{2/0}100,\\
1\overline{1}2&\xrightarrow{0/1}01\overline{1},\ \xrightarrow{1/1}010,\ \xrightarrow{2/1}100,\\
01\overline{1}&\xrightarrow{0/0}001,\ \xrightarrow{1/0}002,\ \xrightarrow{2/0}1\overline{1}2,\\
11\overline{1}&\xrightarrow{0/1}001,\ \xrightarrow{1/1}002,\ \xrightarrow{2/1}1\overline{1}2.
\end{aligned}
\]
Define the terminal output function $\omega:Q\to X_3$ by
\[
\omega(s_1s_2s_3):=\pi_3\!\bigl(Z(s_1F_4+s_2F_3+s_3F_2)\bigr)\in X_3,
\]
where $Z(\cdot)$ is the Zeckendorf digit sequence (Section~\ref{subsec:prelim-zeckendorf}).
For an input word $s\in\{0,1,2\}^m$, the transducer output is the length-$(m+3)$ word obtained by concatenating the edge outputs $e_1\cdots e_m$ with the terminal word $\omega(q_m)$, where $q_m$ is the final state after reading $s$.
\end{definition}

\begin{theorem}[Existence of an online normalizer with delay $3$]\label{thm:online-delay-3}
The subsequential transducer of Definition~\ref{def:delay3-transducer} computes $\Fold_\infty$ on finite-support inputs over $\{0,1,2\}$.
Equivalently, there exists a deterministic finite-state online normalizer for Zeckendorf admissibility with delay $3$.
\end{theorem}
\begin{proof}
This transducer is a standard online normalization device for Fibonacci numeration (see \cite{Frougny1992NumbersAutomata,Frougny1999OnlineAddition}).
We briefly indicate the verification principle.

For each transition $q\xrightarrow{d/e}q'$ one checks that the transformation ``append input digit $d$'' and ``emit output digit $e$'' preserves the Fibonacci value up to a bounded residual represented by the next state.
More precisely, there exists an integer-valued residual functional $R:Q\to\ZZ$ such that, if $q_t$ is the state after reading $t$ digits and $e_t$ is the emitted output at step $t$, then the identity
\[
\sum_{k=1}^{t} d_k F_{k+1}
\;=\;
\sum_{k=1}^{t} e_k F_{k+1}
\;+\;
R(q_t)
\]
holds for all $t$, with $R(q_t)$ depending only on the finite state.
The terminal output $\omega(q_m)$ is exactly the Zeckendorf normalization of the residual value carried by the final state, expressed on the last three weights.
Consequently, the full output (edge outputs plus terminal output) has the same value as the input and is Zeckendorf-admissible (no adjacent $11$), hence it must coincide with the unique normal form $\mathrm{NF}(\iota_m(s))=\Fold_\infty(s)$ by Zeckendorf uniqueness (Theorem~\ref{thm:rewrite-unique-normal-form}).
\end{proof}

\subsection{Optimality: delay-$2$ is impossible}\label{app:online-delay-optimality}

\begin{proposition}[Minimal delay is at least $3$]\label{prop:online-delay-min-3}
Any deterministic online transducer computing $\Fold_\infty$ on $\{0,1,2\}^\NN$ must have delay at least $3$.
\end{proposition}
\begin{proof}
Consider the two finite-support inputs (viewed as configurations in $\mathcal{D}$):
\[
s:=1000,\qquad s':=1002,
\]
so $s$ and $s'$ agree on the first three digits.
Their values are
\[
\mathrm{Val}(\iota_4(s))=F_2=1,
\qquad
\mathrm{Val}(\iota_4(s'))=F_2+2F_5=1+2\cdot 5=11.
\]
The Zeckendorf representation of $1$ has first digit $1$, while the Zeckendorf representation of $11=8+3=F_6+F_4$ has first digit $0$.
Therefore
\[
\bigl(\Fold_\infty(s)\bigr)_1\neq \bigl(\Fold_\infty(s')\bigr)_1,
\]
even though $s$ and $s'$ coincide on the input prefix of length $3$.
Hence no deterministic online algorithm that outputs the first normalized digit after reading at most two additional input digits (delay $\le 2$) can compute $\Fold_\infty$.
\end{proof}

